% ------------------------------------------------------------------
%  Capitolo 3 — Il desiderio di sapere
%  Parte I
%  Ricerca di riferimento: ricerca 01 §8, §10; 03 §2.5
%  Voci bibliografiche principali: recht1988, willingham2007, hulleman2009, yeager2014
%  Epigrafe: ◐ verificata indirettamente (vedi ricerca 03 e registro, ricerca 00)
%  Scheda del capitolo: ricerca/06_indice-ragionato.md, capitolo 3
% ------------------------------------------------------------------
\chapter{Il desiderio di sapere}
\label{cap:sapere}

\epigrafe[Tutti gli uomini per natura desiderano sapere.]{Πάντες ἄνθρωποι τοῦ εἰδέναι ὀρέγονται φύσει.}{Aristotele, Metafisica A 1, 980a21}

\iniziale{Q}{uesto capitolo} \segnaposto{apertura: da scrivere — vedi la scheda nell'indice ragionato}.

\section{Tutti gli uomini desiderano sapere}

\segnaposto{da scrivere}

\section{Si pensa con ciò che si sa}

\segnaposto{da scrivere}

\section{Sapere per essere liberi}

\segnaposto{da scrivere}

\section{Dare un senso a ciò che si studia}

\segnaposto{da scrivere}

\begin{daricordare}
\segnaposto{il principio del capitolo in due righe}
\end{daricordare}

\begin{domande}
  \item \segnaposto{domanda di recupero su questo capitolo}
  \item \segnaposto{domanda di applicazione a una materia che studi}
  \item \segnaposto{domanda di ripresa su un capitolo precedente}
\end{domande}

\begin{sintesi}
  \item \segnaposto{punto di sintesi}
\end{sintesi}

\finecapitolo
